%
% Documento: Resumo (Inglês)
%

\begin{center}
\imprimirtitleenglish
\end{center}

\begin{resumo}[Abstract]
The development of modern front-end applications demands ecosystems that balance delivery agility, maintainability, and code reliability. This study investigates perceptions and factors associated with the adoption of JavaScript (JS), characterized by dynamic typing, and TypeScript (TS), which provides optional static typing, comparing novice and experienced developers. A mixed-methods survey was conducted with a convenience sample of 23 front-end developers. Data were analyzed using descriptive statistics and thematic categorization of open-ended responses. In the sample, novices expressed stronger agreement regarding JavaScript's rapid prototyping (mean 4.29), whereas experienced participants reported greater productivity with TypeScript (4.22, compared with 2.93 among novices) and a stronger preference for TypeScript in new projects (4.56). Participants also associated TypeScript with early bug detection (4.17 overall), maintainability (4.44 among experienced participants), and team standardization (4.30). Qualitative accounts, however, identified configuration costs, migration difficulties, library contracts, and the limits of static checking at runtime. The findings suggest that language choice is contextual and related to project scale, maintenance horizon, and team experience; they do not support generalizations to the entire developer community.

\vspace{\onelineskip}
\noindent
\textbf{Keywords}: JavaScript. TypeScript. Front-End Development. Static Typing. Software Maintainability. Software Engineering.
\end{resumo}
